\exercisesection{\S~2}

\begin{enumerate}
\def\labelenumi{\arabic{enumi})}
\tightlist
\item
  Let \(\mathbf{G}\) be a connected compact Lie group, \(g\) an element of \(\mathbf{G}\) .
\end{enumerate}

\begin{enumerate}
\def\labelenumi{\alph{enumi})}
\item
  Show that there exists an integer \(n \geq 1\) such that the centralizer \(Z(g^{n})\) is connected (prove that, for suitable n, the closed subgroup generated by \(g^{n}\) is a torus).
\item
  Assume that the dimension of \(\operatorname{Ker}(\operatorname{Ad}g^n - 1)\) is independent of \(n\) ( \(n \geq 1\) ). Prove that \(Z(g)\) is connected.
\item
  If \(g^n\) is regular for all \(n \geq 1\) , \(Z(g)\) is connected.
\end{enumerate}

\begin{enumerate}
\def\labelenumi{\arabic{enumi})}
\setcounter{enumi}{1}
\item
  Show that every connected compact Lie group G is the semi-direct product of its derived group by a torus (observe that, if T is a maximal torus of G, then \(D(G) \cap T\) is a torus).
\item
  Let G be a compact Lie group, \(\mathfrak{g}\) its Lie algebra, \(\mathfrak{s}\) a vector subspace of \(\mathfrak{g}\) such that, for \(x, y, z\) in \(\mathfrak{s}\), we have \([x, [y, z]] \in \mathfrak{s}\) . Let \(\mathcal{L}\) be the set of commutative subalgebras of \(\mathfrak{g}\) contained in \(\mathfrak{s}\). Show that the identity component of the stabilizer of \(\mathfrak{s}\) in G operates transitively on the set of maximal elements of \(\mathcal{L}\) (argue as in the proof of Th. 1).
\item
  Let G be a connected compact Lie group, T a maximal torus of G and S a sub-torus of T. Denote by \(\Sigma\) (resp. F) the stabilizer (resp. the fixer) of S in \(W_{G}(T)\) .
\end{enumerate}

\begin{enumerate}
\def\labelenumi{\alph{enumi})}
\item
  Prove that the group \(\mathrm{N_G(S) / Z_G(S)}\) is isomorphic to the quotient \(\Sigma /\mathrm{F}\) .
\item
  Let H be a connected closed subgroup of G, such that S is a maximal torus of H. Show that every element of \(W_{\mathrm{H}}(\mathrm{S})\) , considered as an automorphism of S, is the restriction to S of an element of \(\Sigma\) .
\end{enumerate}

\begin{enumerate}
\def\labelenumi{\arabic{enumi})}
\setcounter{enumi}{4}
\tightlist
\item
  Let \(G\) be a connected compact Lie group and \(S\) a torus of \(G\) . Show that the following conditions are equivalent:
\end{enumerate}

\begin{enumerate}
\def\labelenumi{(\roman{enumi})}
\item
  S is contained in a unique maximal torus;
\item
  \(Z_{G}(S)\) is a maximal torus;
\item
  S contains a regular element.
\end{enumerate}

\begin{enumerate}
\def\labelenumi{\arabic{enumi})}
\setcounter{enumi}{5}
\tightlist
\item
  Let G be a connected compact Lie group, T a maximal torus of G, \(\mathfrak{g}\) (resp. \(\mathfrak{t}\)) the Lie algebra of G (resp. T), and \(i : \mathfrak{t} \to \mathfrak{g}\) the canonical injection. Prove that the map \(^{t}i : \mathfrak{g}^{*} \to \mathfrak{t}^{*}\) induces by passage to the quotient a homeomorphism from \(\mathfrak{g}^{*}/G\) to \(\mathfrak{t}^{*}/W_{G}(T)\) .
\end{enumerate}

¶ 7) Let X and Y be two real manifolds of class \(C^{r}\) ( \(1 \leq r \leq \omega\) ), separated, connected and of dimension n; let \(f : X \to Y\) be a proper morphism of class \(C^{r}\) , equipped with an orientation F (Differentiable and Analytic Manifolds, Results, 10.2.5).

\begin{enumerate}
\def\labelenumi{\alph{enumi})}
\item
  Show that there exists a unique real number d such that \(\int_{X} f^{*}\alpha = d\int_{Y}\alpha\) for every twisted differential form \(\alpha\) of degree n on Y of class \(C^{1}\) and of compact support (use Differentiable and Analytic Manifolds, Results, 11.2.4).
\item
  Let \(y \in \mathbf{Y}\) be such that \(f\) is étale at every point of \(f^{-1}(y)\) . For \(x \in f^{-1}(y)\) , put \(\nu_x(f) = 1\) (resp. \(\nu_x(f) = -1\) ) if the maps \(\mathrm{F}_x\) and \(\tilde{f}_x\) (Differentiable and Analytic Manifolds, Results, 10.2.5, Example b)) from \(\mathrm{Or}(\mathrm{T}_x(\mathrm{X}))\) to \(\mathrm{Or}(\mathrm{T}_y(\mathrm{Y}))\) coincide (resp. are opposite). Prove that \(d = \sum_{x \in f^{-1}(y)} \nu_x(f)\) and, in particular, that \(d \in \mathbf{Z}\) .
\end{enumerate}

We call d the degree of f and denote it by \(\deg f\) . If X = Y and X is orientable, we can take for F the orientation that preserves the orientation of X.

\begin{enumerate}
\def\labelenumi{\alph{enumi})}
\setcounter{enumi}{2}
\item
  Show that, if \(\deg f \neq 0\) then \(f\) is surjective.
\item
  If there exists \(x \in X\) such that \(f^{-1}(f(x)) = \{x\}\) and such that f is étale at x, then f is surjective.
\end{enumerate}

\begin{enumerate}
\def\labelenumi{\arabic{enumi})}
\setcounter{enumi}{7}
\tightlist
\item
  Let G be a connected compact Lie group, dg the normalized Haar measure on G.
\end{enumerate}

\begin{enumerate}
\def\labelenumi{\alph{enumi})}
\tightlist
\item
  Let \(f: \mathrm{G} \to \mathrm{G}\) be a morphism of manifolds; for \(g \in \mathrm{G}\) , the differential \(\mathrm{T}_g(f)\) can be identified by left translations with a linear map from \(\mathrm{L}(\mathrm{G})\) to \(\mathrm{L}(\mathrm{G})\) . Prove the formulas
\end{enumerate}

\[
\deg f. \int_ {\mathrm{G}} \varphi (g) d g = \int_ {\mathrm{G}} \varphi (f (g)) \det \mathrm{T} _ {g} (f) d g
\]

for every integrable (complex-valued) function \(\varphi\) on G, and

\[
\deg f = \int_ {\mathrm{G}} \det \mathrm{T} _ {g} (f) d g.
\]

\begin{enumerate}
\def\labelenumi{\alph{enumi})}
\setcounter{enumi}{1}
\tightlist
\item
  Let \(\psi_k: \mathbf{G} \to \mathbf{G}\) be the map \(g \mapsto g^k\) . Prove the formula
\end{enumerate}

\[
\deg \psi_ {k} = \int_ {\mathrm{G}} \det (1 + \operatorname{Ad} g + \dots + (\operatorname{Ad} g) ^ {k - 1}) d g.
\]

\begin{enumerate}
\def\labelenumi{\alph{enumi})}
\setcounter{enumi}{2}
\tightlist
\item
  Let \(\mathrm{T}\) be a maximal torus of \(\mathrm{G}\) , and \(\mathrm{T}_r\) the set of regular elements of \(\mathrm{T}\) . Show that \(\exp_{\mathrm{G}}^{-1}(\mathrm{T}_r) \subset \mathrm{L}(\mathrm{T})\) and \(\psi_k^{-1}(\mathrm{T}_r) \subset \mathrm{T}_r\) for all \(k \geq 1\) .
\end{enumerate}

\(d\)) Show that \(\deg \psi_{k} = k^{\dim (\mathrm{T})}\) ; deduce the equality

\[
\int_ {\mathrm{G}} \det (1 + \operatorname{Ad} g + \dots + (\operatorname{Ad} g) ^ {k - 1}) d g = k ^ {\dim (\mathrm{T})}.
\]

\begin{enumerate}
\def\labelenumi{\arabic{enumi})}
\setcounter{enumi}{8}
\tightlist
\item
  This exercise is devoted to another proof of Th. 2. Let G be a connected compact Lie group, T a maximal torus of G, N its normalizer. Denote by \(G \times^{N} T\) the quotient manifold of \(G \times T\) by N for the operation defined by \((g, t).n = (gn, n^{-1}tn)\) (Differentiable and Analytic Manifolds, Results, 6.5.1).
\end{enumerate}

\begin{enumerate}
\def\labelenumi{\alph{enumi})}
\item
  Show that the morphism \((g, t) \mapsto gtg^{-1}\) from \(G \times T\) to G defines by passage to the quotient an analytic morphism \(f : G \times^{N} T \to G\) .
\item
  Let \(\theta\) be an element of T whose set of powers is dense in T (General Topology, Chap. VII, §1, no. 3, Cor. 2 of Prop. 7). Show that \(f^{-1}(\theta) = \{x\}\) , where \(x\) is the class of \((e,\theta)\) in \(\mathrm{G} \times^{\mathrm{N}}\mathrm{T}\) , and that \(f\) is étale at \(x\) .
\item
  Conclude from Exerc. 7 d) that \(f\) is surjective, and deduce another proof of Th. 2.
\end{enumerate}

\begin{enumerate}
\def\labelenumi{\arabic{enumi})}
\setcounter{enumi}{9}
\tightlist
\item
  * In this exercise, we use the following result from algebraic topology (Lefschetz's formula \(^{11}\) ): let X be a finite dimensional compact manifold, and f a morphism from X to itself. Assume that the set F of points x of X such that \(f(x) = x\) is finite, and that for all \(x \in F\) the number \(\delta(x) = \det(1 - \mathrm{T}_{x}(f))\) is non-zero. If \(\mathrm{H}^{i}(f)\) denotes the endomorphism of the vector space \(\mathrm{H}^{i}(\mathrm{X}, \mathbf{R})\) induced by f (for \(i \geq 0\) ), then
\end{enumerate}

\[
\sum_ {x \in \mathrm{F}} \delta (x) / | \delta (x) | = \sum_ {i \geq 0} (- 1) ^ {i} \operatorname{Tr} \mathrm{H} ^ {i} (f).
\]

Let G be a connected compact Lie group, T a maximal torus of G. For \(g \in G\) , denote by \(\tau(g)\) the automorphism of the manifold G/T induced by left multiplication by g.

\begin{enumerate}
\def\labelenumi{\alph{enumi})}
\item
  Let \(t\) be an element of \(\mathrm{T}\) such that the subgroup generated by \(t\) is dense in \(\mathrm{T}\) . Show that the fixed points of \(\tau(t)\) are the classes \(n\mathrm{T}\) for \(n \in \mathrm{N}_{\mathrm{G}}(\mathrm{T})\) . Deduce that \((\mathrm{N}_{\mathrm{G}}(\mathrm{T}):\mathrm{T}) = \sum_{i}(-1)^{i}\dim_{\mathbf{R}}\mathrm{H}^{i}(\mathrm{G / T},\mathbf{R})\) .
\item
  Let g be an arbitrary element of G; show that the set of fixed points of \(\tau(g)\) is non-empty, and deduce another proof of Th. 2.*
\end{enumerate}

\begin{enumerate}
\def\labelenumi{\arabic{enumi})}
\setcounter{enumi}{10}
\tightlist
\item
  Let G be a compact Lie group. A subgroup S of G is said to be of type (C) if it is equal to the closure of a cyclic subgroup that is of finite index in its normalizer.
\end{enumerate}

\begin{enumerate}
\def\labelenumi{\alph{enumi})}
\item
  Let \(S\) be a subgroup of type (C); show that \(S_0\) is a maximal torus of \(G_0\) , and that \(S\) is the direct product of \(S_0\) and a finite cyclic subgroup.
\item
  Prove that every element g of G is contained in a subgroup of type (C) (consider the group generated by g and a maximal torus of \(Z(g)_{0}\) ).
\item
  Let S be a subgroup of type (C), and s an element of S whose set of powers is dense in S. Show that every element of \(sG_{0}\) is conjugate under \(\operatorname{Int}(G_{0})\) to an element of \(sS_{0}\) (use the method of Exerc. 10).
\item
  Denote by \(p: G \to G/G_{0}\) the quotient map. Show that the map \(S \mapsto p(S)\) induces a bijection from the set of conjugacy classes of subgroups of G of type (C) to the set of conjugacy classes of cyclic subgroups of \(G/G_{0}\) .
\item
  Let S be a subgroup of G of type (C); denote by \(S_{\rho}\) the set of elements of S whose class in \(S/S_{0}\) generates \(S/S_{0}\) . Show that two elements of \(S_{\rho}\) that are conjugate in G are conjugate in \(N_{G}(S)\) .
\end{enumerate}
% semantic-fragments: legacy-input-seam
\relax{}
